\section{Prospective fourth-order response prediction}\label{sec:quartic}

\subsection{Symmetric local response}

Introduce normalized error coordinates $z\in\Rr^{3}$ through
\[
\eps=Dz,\qquad D=\diag(0.06,0.04,0.05).
\]
For a direction $v\in\Rr^{3}$, define the symmetric response
\begin{equation}\label{eq:symmetric}
S_{\gamma}(t,v)=\frac{\Iloc_{\gamma}(tDv)+\Iloc_{\gamma}(-tDv)}{2}.
\end{equation}
Analyticity of the finite matrix exponential gives
\begin{equation}\label{eq:taylor}
S_{\gamma}(t,v)=q_{2,\gamma}(v)t^{2}+q_{4,\gamma}(v)t^{4}+q_{6,\gamma}(v)t^{6}+\cdots.
\end{equation}
Odd powers cancel by construction.

By Lemma~\ref{lem:quadratic}, exact projective first-response matching forces the quadratic
loss forms of matched controls to coincide exactly, so the quartic term is the first
possible symmetric discriminator. The numerical cohorts here are matched only within
tolerance; consistently, the relative spread of the averaged quadratic coefficient across
the validation cohorts is approximately $10^{-7}$. The first useful discriminator is
therefore the quartic coefficient.

\subsection{Prospective quartic ranking}

The quartic rule developed here is a \emph{geometric predictor}, not a universal theorem:
it is gap-dependent and certifies only a subset of pairwise comparisons, a boundary made
precise in Sections~\ref{sec:covariant}--\ref{sec:certification}.

\subsubsection{Discovery and validation separation}

The initial four named controls were used only for discovery and orientation. A twelve-path
cohort was then used to compare candidate formulas. The smallest sufficient rule on that
training cohort was
\[
\text{smaller }\Gg_{4}\ \Longrightarrow\ \text{smaller mean finite-error infidelity.}
\]
The addition of the sixth-order coefficient did not improve the training rank correlation,
so the quartic rule was frozen.

For the prospective audit, the coefficients $q_{2},q_{4},q_{6}$ were estimated by a frozen
least-squares fit of
\[
S_{\gamma}(t,e_{a})=q_{2,a}t^{2}+q_{4,a}t^{4}+q_{6,a}t^{6}
\]
at
\[
t\in\Bigl\{\tfrac{1}{32},\tfrac{1}{24},\tfrac{1}{20},\tfrac{1}{16},
\tfrac{1}{14},\tfrac{1}{12},\tfrac{1}{10},\tfrac{1}{8}\Bigr\}.
\]
The tensor audit below is computationally distinct: its directional derivatives are
evaluated directly at $t=0$ by block-matrix exponentials, rather than by fitting
nonzero-radius samples.

A new twenty-path cohort was generated with a different seed. For each candidate, a ranking
certificate containing phases, local coefficients, and predicted order was written and
SHA-256 hashed before the held-out outcome evaluator was enabled. The predeclared primary
gate was
\[
\rho_{\mathrm{Spearman}}\ge 0.95,
\]
together with top-path agreement within each recorded execution.

\subsubsection{Quartic validation result}

The prospective validation result is therefore reported at the predeclared threshold level:
\begin{equation}\label{eq:spearman}
\rho_{\mathrm{Spearman}}\bigl(-\Gg_{4},-\II\bigr)\ge0.95.
\end{equation}
Across the recorded architecture-specific executions ($n=2$) of the prospective
cohort-generation-and-evaluation protocol, the predeclared rank gate was satisfied. The
arm64 reference environment recorded $\rho=0.99398$ with top path pv17, and the
x86\_64/Rosetta diagnostic environment recorded $\rho=0.98947$ with top path pv02; thus the
exact correlation, complete ranking, and named top path were not architecture-invariant,
while both environments cleared the $\rho\ge0.95$ gate. This empirical evidence is not
used in the proof of the exact-root ordering theorem and does not establish a universal
fourth-order ranking law. Because the quartic rule, ranking direction, acceptance threshold,
and ranking certificate were frozen before held-out finite-error evaluation, the
threshold-level result is not an in-sample fit; nevertheless, it remains cohort-specific
prospective evidence rather than a theorem. Two local coefficient-extraction windows had
rank correlation one, and the maximum relative fit residual was $2.37\times10^{-10}$, the
direction-safe upward rendering of the frozen field
\texttt{maximum\_relative\_fit\_residual} in \texttt{results/g4\_prospective/report.json}.

Subsequent cohorts show that $\Gg_{4}$ is not universally sufficient: depending on the
generated matched paths, the prospective quartic correlation varies appreciably across
seeded cohorts. This variation motivated the bounded higher-order analysis rather than post
hoc replacement of the quartic formula.

\section{Covariant quartic response geometry}\label{sec:covariant}

\subsection{Quartic tensor}

The object recovered in this section is a quartic susceptibility of the endpoint loss---a
component of the fourth jet of the endpoint map composed with infidelity---not a
Riemannian curvature or an instance of the quantum geometric tensor~\cite{provost1980}.
There exists a symmetric fourth-order tensor $\Aa^{(4)}_{\gamma}$ such that
\begin{equation}\label{eq:quartictensor}
q_{4,\gamma}(v)=\Aa^{(4)}_{\gamma,ijkl}\,v_{i}v_{j}v_{k}v_{l}.
\end{equation}
In three dimensions a symmetric fourth-order tensor has $\binom{3+4-1}{4}=15$ independent
components. The tensor $\Aa^{(4)}_{\gamma}$ is reconstructed from the fourth jet of the
path-local loss $\Iloc_{\gamma}$: we recover its components from 36 seeded directions,
whose directional Taylor coefficients are computed directly at $t=0$ through order six by
block-matrix exponentials, rather than by fitting nonzero-radius samples.

Let
\[
z\sim\mathrm{Unif}\{\pm e_{\Omega},\pm e_{\Delta},\pm e_{V}\},\qquad \eps=Dz,
\]
so that the physical error scales $0.06^{4},0.04^{4},0.05^{4}$ are carried by $D$ and
absorbed into the tensor components. For the six-axis law in Eq.~\eqref{eq:errorlaw},
define the fourth moment
\[
M^{(4)}_{ijkl}=\ev{z_{i}z_{j}z_{k}z_{l}}.
\]
The scalar quartic predictor is
\begin{equation}\label{eq:g4}
\Gg_{4}(\gamma;M^{(4)})=\Aa^{(4)}_{\gamma,ijkl}M^{(4)}_{ijkl}.
\end{equation}
For the equally weighted signed axes,
\[
\Gg_{4}=\frac{1}{3}\Bigl(\Aa^{(4)}_{\Omega\Omega\Omega\Omega}
+\Aa^{(4)}_{\Delta\Delta\Delta\Delta}+\Aa^{(4)}_{VVVV}\Bigr).
\]

\subsection{Coordinate covariance}

\begin{proposition}[Covariance and invariant contraction]\label{prop:covariance}
Let $z=Ry$ for an invertible real matrix $R$. Then
\begin{equation}\label{eq:atransform}
\Aa^{(4)}_{y,abcd}=\Aa^{(4)}_{z,ijkl}\,R_{ia}R_{jb}R_{kc}R_{ld},
\end{equation}
\begin{equation}\label{eq:mtransform}
M^{(4)}_{y,abcd}=(R^{-1})_{ai}(R^{-1})_{bj}(R^{-1})_{ck}(R^{-1})_{dl}\,M^{(4)}_{z,ijkl},
\end{equation}
and
\begin{equation}\label{eq:invariant}
\Aa^{(4)}_{y}:M^{(4)}_{y}=\Aa^{(4)}_{z}:M^{(4)}_{z}.
\end{equation}
\end{proposition}

\begin{proof}
Equation~\eqref{eq:atransform} follows by substituting $z=Ry$ into the quartic form.
Equation~\eqref{eq:mtransform} follows from $y=R^{-1}z$. Contracting all four indices
cancels each $R$ with the corresponding $R^{-1}$, proving Eq.~\eqref{eq:invariant}.
\end{proof}

This proposition is algebraic. The numerical audit additionally recomputes the tensor
directly in the transformed coordinates. For a nonorthogonal map of condition number
approximately $1.82$, the direct-refit relative error is $1.08\times10^{-14}$ and the
maximum invariant-contraction relative error is $2.36\times10^{-15}$.

The covariant object is not the scalar $\Gg_{4}$ in isolation. It is the fourth-order
response tensor $\Aa^{(4)}$ together with the physical fourth noise moment $M^{(4)}$. Their
contraction is independent of a linear reparameterization of error coordinates. This avoids
interpreting a coordinate-specific axis average as an intrinsic curvature.

Establishing a relation of this susceptibility to a canonical connection or curvature on a
control quotient would require additional structure.

For a noise moment $M$, the low-cost screening problem is
\[
\gamma^{\star}\in\argmin_{\gamma\in\Ff_{\eta_{\mathrm{match}}}}\ \Aa^{(4)}_{\gamma}:M^{(4)}.
\]
Pairs that fail the quartic gap test can be escalated adaptively to sixth, eighth, or
higher jet order until their intervals separate. This suggests a certified branch-and-bound
control selector: compute low-order local tensors for all candidates, refine only ambiguous
pairs, and reserve direct finite-error simulation or hardware evaluation for model
validation rather than exhaustive ranking.

The gap-dependent certification criterion used by the implementation retains the residual
constant and quadratic terms explicitly. The certified quantity is the six-axis averaged
reference loss at radius $t$,
\begin{equation}\label{eq:radiusloss}
\Iref_{k}(t)=\frac{1}{6}\sum_{a\in\{\Omega,\Delta,V\}}\sum_{\sigma=\pm1}
\Iref_{k}(\sigma t D e_{a}).
\end{equation}
At unit radius $t=1$ the six signed probe points $\sigma D e_{a}$ are exactly the six
signed axial errors of the declared error set, so Eq.~\eqref{eq:radiusloss} reduces to the
six-axis mean functional of Eq.~\eqref{eq:meanloss}; the radius-$t$ form is its
one-parameter rescaling used for the local jet analysis.
For each path $k$ write the truncated local model and its enclosure as
\begin{equation}\label{eq:truncated}
P_{k}(t)=C_{0,k}+C_{2,k}t^{2}+\Gg_{4,k}t^{4},\qquad
\Iref_{k}(t)\in P_{k}(t)+[-B_{k}(t),\,B_{k}(t)].
\end{equation}
If
\begin{equation}\label{eq:gaptest}
\lvert P_{a}(t)-P_{b}(t)\rvert > B_{a}(t)+B_{b}(t),
\end{equation}
then the sign of $\Iref_{a}(t)-\Iref_{b}(t)$ is strictly certified. Only when
$C_{0,a}=C_{0,b}$ and $C_{2,a}=C_{2,b}$ does this reduce to the simplified quartic-gap form
\[
\lvert\Gg_{4,a}-\Gg_{4,b}\rvert\,\lvert t\rvert^{4} > B_{a}(t)+B_{b}(t).
\]
The implementation folds the residual $C_{0}$ and $C_{2}$ differences into the correction
intervals, so they are not omitted from the certificate.

\section{Computer-assisted certification}\label{sec:certification}

\subsection{Krawczyk root theorem}

The numerical paths satisfy the endpoint and first-order matching constraints to order
$10^{-11}$ (see Remark~\ref{rem:residual} for the precise norm and its relation to the
formal boxes). To move from numerical candidates to exact mathematical roots, we use the
projective constraint map $F:\Rr^{24}\to\Rr^{16}$ of Eq.~\eqref{eq:constraintmap}. We do
\emph{not} split the 24 original phase coordinates into a fixed subset of 16 transverse
coordinates and 8 retained fibre coordinates. Instead, for each frozen path $k$ the
certificate uses the affine transverse parameterization
\begin{equation}\label{eq:affine}
\gamma=\widehat{\gamma}_{k}+N_{k}\,x,
\qquad x\in X=[-3\times10^{-12},\,3\times10^{-12}]^{16}\subset\Rr^{16},
\end{equation}
where $\widehat{\gamma}_{k}\in\Rr^{24}$ is the frozen box centre of path $k$ and the
columns of the frozen $24\times16$ matrix $N_{k}$ form the declared orthonormal transverse
basis. The variable $x$ is the vector of transverse-basis coefficients, not a selection of
16 of the original 24 phase coordinates, and $N_{k}$ varies from path to path. On this
declared affine transverse slice we consider the reduced square map
\[
F_{k}(x)=F(\widehat{\gamma}_{k}+N_{k}x),\qquad F_{k}:\Rr^{16}\to\Rr^{16},
\]
whose Jacobian is the reduced $16\times16$ matrix
\begin{equation}\label{eq:reducedjac}
DF_{k}(x)=DF(\widehat{\gamma}_{k}+N_{k}x)\,N_{k}\in\Rr^{16\times16}.
\end{equation}
All interval operators below act on this path-dependent reduced map; at no point is a
transverse vector $x\in\Rr^{16}$ passed directly to $F:\Rr^{24}\to\Rr^{16}$. For a box
$X\subset\Rr^{16}$ with centre $x_{0}$, the Krawczyk operator for path $k$ is
\begin{equation}\label{eq:krawczyk}
K_{k}(X)=x_{0}-C_{k}F_{k}(x_{0})+\bigl(I-C_{k}[DF_{k}(X)]\bigr)(X-x_{0}),
\end{equation}
where $C_{k}$ is the frozen point preconditioner of the corresponding reduced midpoint
Jacobian $DF_{k}(x_{0})$ (a numerical approximate inverse) and $[DF_{k}(X)]$ is an
interval enclosure of the reduced Jacobian~\eqref{eq:reducedjac} over $X$. In centered
coordinates $x_{0}=0$, as used by the implementation, this reduces to
\[
K_{k}(X)=-C_{k}F_{k}(0)+\bigl(I-C_{k}[DF_{k}(X)]\bigr)X.
\]
The strict inclusion $K_{k}(X)\subset\intt(X)$ then implies that $F_{k}$ has a unique
zero inside $X$~\cite{moore2009,tucker2011}. The projective chart denominator $\psi_{0}$ of
Eq.~\eqref{eq:chart} is formally verified to exclude zero on every accepted box, ensuring
the coordinate chart is valid throughout the Krawczyk computation.

\begin{theorem}[Krawczyk inclusions for 12 declared controls]\label{thm:krawczyk}
For each of the twelve declared controls, with the affine transverse parameterization of
Eq.~\eqref{eq:affine}, the common transverse box
\[
X_{k}=X=[-3\times10^{-12},\,3\times10^{-12}]^{16}
\]
satisfies the strict Krawczyk inclusion
\[
K_{k}(X_{k})\subset\intt(X_{k}).
\]
Consequently, each box $X_{k}$ contains a unique exact root $x^{\ast}_{k}$ of
$F_{k}(x)=0$, i.e., a control
$\gamma^{\ast}_{k}=\widehat{\gamma}_{k}+N_{k}x^{\ast}_{k}$ that is exactly matched in
projective output state and first projective response to the reference. Existence and
uniqueness are certified within the declared transverse box only, not on the full
24-dimensional phase space or the whole response fibre.
\end{theorem}

\begin{proof}[Computer-assisted proof]
A numerical approximate inverse of the midpoint Jacobian is used as the point
preconditioner. The Jacobian over each full box is then enclosed with 192-bit Arb
arithmetic, and the resulting Krawczyk image is verified to lie strictly inside the box.
The projective chart denominator is formally verified to exclude zero on every accepted
box. For all twelve boxes, the inclusion $K_{k}(X)\subset\intt(X)$ is verified; the worst
box-normalized utilization, read from the frozen certificate, is
\[
\max_{k}\frac{\lVert K_{k}(X)\rVert_{\infty}}{3\times10^{-12}}<0.866<1,
\]
so every coordinate image clears the box wall by a factor of at least $1.15$. The SHA-256
protocol and certificate identifiers of this run are recorded in
Appendix~\ref{app:hashes}. By the Krawczyk
theorem~\cite{moore2009,tucker2011,tucker2002}, each box contains a locally unique root.
As a \emph{supplementary} audit-closure check---not an additional premise of the strict
inclusion just established---the regularity of each frozen point preconditioner is
independently certified by a $256$-bit outward-rounded Rump--Neumann argument: for all
twelve paths $\lVert I-R_{k}C_{k}\rVert_{\infty}<1$, so each $C_{k}$ is nonsingular
(certificate \nolinkurl{results/audit_closure/p0_preconditioner_certificate.json}). This
P0 certificate belongs to the later v0.3.2 audit-closure supplement
(Remark~\ref{rem:p0role}); it strengthens auditability but neither replaces nor alters the
v0.3.1 frozen Krawczyk certificate that proves the theorem. In the strict-inclusion form
of the Krawczyk theorem used here, the required regularity is already implied by the
accepted inclusion; the v0.3.2 P0 certificate is therefore an independent redundant audit
of the production preconditioner, not an additional premise of Theorem~\ref{thm:krawczyk}.
\end{proof}

\begin{remark}[Residual, norm, and box relation]\label{rem:residual}
Two distinct residuals must not be conflated. The \emph{unpreconditioned} constraint
residual $\lVert F\rVert$ of the numerical candidates is of order $10^{-11}$; it is a
float64 diagnostic of the candidate-generation optimizer and is \emph{not} compared
directly with the $3\times10^{-12}$ box radius. The strict Krawczyk inclusion is decided
instead by the full interval image
$-C_{k}F_{k}(x_{0})+\bigl(I-C_{k}\,[DF_{k}(X)]\bigr)(X-x_{0})$
of Eq.~\eqref{eq:krawczyk}, whose outward-rounded endpoints are enclosed on the formal
cohort box centres $\widehat{\gamma}_{k}$. The preconditioned image
$\lVert K_{k}(X)\rVert_{\infty}$, already preconditioned by $C_{k}$, is what the
utilization bound $0.866$ above reports. The frozen certificate does not tabulate a single
scalar preconditioned centre residual; accordingly we do not quote one, and the
$10^{-11}$ figure refers only to the unpreconditioned optimizer residual.
\end{remark}

\begin{remark}[Scope of the Krawczyk inclusion theorem]
The theorem establishes local uniqueness inside each declared transverse box. It does not
certify global uniqueness on the entire fibre, nor does it prove that every point on the
numerical fibre has an exact representative. The global fibre structure remains open.
\end{remark}

\subsection{Direct finite-error ordering theorem}

The following theorem is the main result of this work. Its proof uses only the certified
Krawczyk boxes and direct outward-rounded propagation of the finite-error functional. No
Taylor truncation, Cauchy extraction, alias enclosure, or local-jet reconstruction is used.

\begin{theorem}[Exact-root finite-error ordering]\label{thm:ordering}
For each of the twelve declared controls, let $X_{k}$ be the transverse phase box
satisfying a strict Krawczyk inclusion from Theorem~\ref{thm:krawczyk}. Direct
outward-rounded Arb propagation of $X_{k}$ through the declared six-point finite-error
functional produces twelve performance intervals $B_{k}$ covering the complete root box,
\[
B_{k}\supseteq
\Bigl\{\,\overline{\mathcal{E}}\bigl(\widehat{\gamma}_{k}+N_{k}x\bigr):
x\in[-3\times10^{-12},\,3\times10^{-12}]^{16}\Bigr\},
\]
that are pairwise disjoint and totally
ordered according to the pre-outcome ordering frozen for the independent twelve-path
formal cohort. Therefore the locally
unique roots exactly matched in projective output state and first projective response
inside the boxes obey the complete frozen finite-error ordering.
\end{theorem}

\begin{proof}[Computer-assisted proof]
By Theorem~\ref{thm:krawczyk}, the locally unique projectively matched control
$\gamma^{\ast}_{k}=\widehat{\gamma}_{k}+N_{k}x^{\ast}_{k}$ belongs to the certified box.
For every box $X_{k}$, the six declared finite-error Hamiltonians are evaluated directly
using 192-bit outward-rounded Arb/ACB arithmetic over the \emph{whole} declared transverse
box, not merely at its centre or at the approximate root: every phase-dependent
trigonometric factor, Hamiltonian matrix exponential, state propagation, overlap,
infidelity, and six-point mean is enclosed. Denote the resulting real performance interval
by $B_{k}$. Because the propagation is over the full box,
\[
B_{k}\supseteq
\Bigl\{\,\overline{\mathcal{E}}\bigl(\widehat{\gamma}_{k}+N_{k}x\bigr):
x\in[-3\times10^{-12},\,3\times10^{-12}]^{16}\Bigr\},
\]
where $\overline{\mathcal{E}}$ is the six-axis mean finite-error functional of
Eq.~\eqref{eq:meanloss}. Since the certified root satisfies $x^{\ast}_{k}\in X$, its
exact-root performance lies in the enclosure,
\[
\II(\gamma^{\ast}_{k})\in B_{k}.
\]
No Taylor truncation, order-30 jet, or local-jet reconstruction enters $B_{k}$. For every
frozen ordered pair $a\prec b$, direct interval comparison verifies
\[
\sup B_{a}<\inf B_{b}.
\]
All $\binom{12}{2}=66$ comparisons pass and have the orientation specified by the frozen
ordering certificate. Because each $B_{k}$ encloses its exact root's performance, the
strict box separation $\sup B_{a}<\inf B_{b}$ implies the strict exact-root ordering
\[
\II(\gamma^{\ast}_{a})<\II(\gamma^{\ast}_{b})
\]
for every frozen ordered pair.
\end{proof}

The predecessor order was frozen before direct propagation; the resulting protocol and
certificate hashes are reported in Appendix~\ref{app:hashes}.

\begin{remark}[Scope of Theorem~\ref{thm:ordering}]
The theorem is conditional on the serialized-decimal finite Hamiltonian model, the declared
transverse boxes, and the six-axis mean error functional. It does not certify model
discrepancy, global fibre topology, many-body scaling, or hardware. It also does not prove
a global eight-dimensional exact fibre theorem; it proves finite-error ordering for the
locally unique exact roots inside the declared boxes.
\end{remark}

\begin{corollary}[Ordering stability under bounded loss discrepancy]\label{cor:stability}
Let $B_{k}$ denote the certified direct performance interval for the $k$th exact root, and
define the minimum certified ordering margin
\[
\Delta_{\min}=\min_{a\prec b}\bigl(\inf B_{b}-\sup B_{a}\bigr)>0,
\]
where $a\prec b$ ranges over the frozen ordered path pairs. Suppose an additional physical
effect changes the evaluated loss of every fixed certified control by at most $\eta$:
\[
\bigl\lvert\widetilde{\II}(\gamma^{\ast}_{k})-\II(\gamma^{\ast}_{k})\bigr\rvert\le\eta
\quad\text{for every }k.
\]
If
\[
2\eta<\Delta_{\min},
\]
then the complete frozen ordering of the twelve fixed certified controls is preserved under
the perturbed loss.
\end{corollary}

\begin{proof}
For every frozen ordered pair $a\prec b$,
\[
\widetilde{\II}(\gamma^{\ast}_{a})\le\sup B_{a}+\eta
<\inf B_{b}-\eta\le\widetilde{\II}(\gamma^{\ast}_{b}),
\]
where the strict middle inequality follows from $2\eta<\Delta_{\min}$.
\end{proof}

For the frozen twelve-root certificate, the minimum formally certified pairwise margin,
computed from the outward-rounded endpoints of the frozen direct-certificate intervals, is
\[
\Delta_{\min}=2.50\times10^{-5},
\]
attained by the pair pv08 and pv11.
For the certified minimum separation $\Delta_{\min}=2.505874\ldots\times10^{-5}$, the
sufficient uniform perturbation condition in Corollary~\ref{cor:stability} becomes
$\eta<\Delta_{\min}/2\approx1.25\times10^{-5}$. This is a stringent model-discrepancy
requirement and should not be interpreted as a certificate of hardware-level ordering
robustness.

Corollary~\ref{cor:stability} is a conditional stability statement for the losses evaluated
at the same certified controls. It does not prove that an added Lindblad term, waveform
filter, calibration drift, or other model perturbation satisfies the required bound, nor
that the controls remain exact response-matched roots of the perturbed model. Proving root
continuation under model perturbation would require a parameterized Krawczyk or validated
implicit-function argument.

\subsection{Order-thirty mechanism certificate}

The following proposition provides a separate, computationally distinct mechanism
certificate using the zero-error jet. Unlike the main theorem, it uses Taylor coefficients,
Cauchy extraction, alias enclosures, and Dyson tail bounds. It is not part of the main
theorem proof.

The floating-point precursor obtains the order-thirty jet by block-matrix exponentials. The
formal phase-box certificate used in Proposition~\ref{prop:order30} instead extracts the
coefficients with a 64-point outward-rounded Cauchy transform.

\begin{proposition}[Order-thirty mechanism certificate]\label{prop:order30}
Propagation of the zero-error jet through order thirty over the twelve Krawczyk phase
boxes, including the formal Cauchy-alias and analytic order-32 tail enclosures, certifies
52/66 frozen path pairs. Every certified orientation agrees with the frozen order, and no
reversed pair is certified.
\end{proposition}

\begin{proof}[Computer-assisted proof]
For path $k$ and physical error axis $a$, define the analytic transition amplitude
\[
g_{k,a}(\lambda)=\braket{\psi_{\mathrm{ref}}}{\psi_{\gamma_{k}}(\lambda De_{a})}
=\sum_{n\ge0}g_{k,a,n}\,\lambda^{n}.
\]
Its coefficients are extracted by a 64-point outward-rounded Cauchy discrete Fourier
transform on the complex unit circle:
\begin{equation}\label{eq:cauchy}
\widehat{g}_{n}=\frac{1}{64}\sum_{k=0}^{63}g(\zeta_{k})\,\zeta_{k}^{-n},
\qquad \zeta_{k}=e^{2\pi i k/64}.
\end{equation}
The discrete transform aliases coefficients $g_{n+64\ell}$. A path-uniform Dyson
bound~\cite{dyson1949} encloses this alias. The implemented enclosure is $10^{-40}$, and the formal gate verifies
that
\[
\frac{e^{K_{a}}K_{a}^{64}}{64!}<10^{-40}
\]
for every physical error axis, where the dependence is parameterized as
\[
H_{a}(s;\lambda_{a})=H_{0}(s)+\lambda_{a}H_{1,a}(s),
\]
with $s$ the physical time and $\lambda_{a}$ the dimensionless error parameter along axis
$a$, and
\[
K_{a}=\int_{0}^{T}\lVert H_{1,a}(s)\rVert_{2}\,ds
=\int_{0}^{T}\left\lVert
\frac{\partial H_{a}(s;\lambda_{a})}{\partial\lambda_{a}}
\bigg|_{\lambda_{a}=0}\right\rVert_{2}ds
\]
is the integrated spectral norm of the Hamiltonian perturbation along axis $a$. Fidelity
coefficients are formed by ball convolution,
\[
F_{n}=\sum_{r=0}^{n}g_{r}\,\bar{g}_{n-r},
\]
and the infidelity coefficients follow from $1-F$. The mean six-axis analytic even Dyson
tail after order thirty is enclosed by $2\times10^{-11}$; the largest individual-axis tail
is below $4\times10^{-11}$.

Propagating the resulting order-thirty jet over the twelve Krawczyk phase boxes yields
performance intervals. Pairwise comparison of the resulting order-thirty intervals
certifies 52/66 frozen path pairs. Every certified orientation agrees with the frozen
order, and no reversed pair is certified. The unresolved pairs arise from dependency and
wrapping in the phase-box propagation of the jet enclosure. They are not reversed physical
comparisons and do not indicate an omitted-tail failure.
\end{proof}
